\title{MusiteDeep: a deep-learning based webserver for protein post-translational modification site prediction and visualization}
\begin{document}
\maketitle

\begin{abstract}
MusiteDeep is an online resource providing a deep- learning framework for protein post-translational modification (PTM) site prediction and visualization. The predictor only uses protein sequences as input and no complex features are needed, which results in a real-time prediction for a large number of pro- teins. It takes less than three minutes to predict for 1000 sequences per PTM type. The output is pre- sented at the amino acid level for the user-selected PTM types. The framework has been benchmarked and has demonstrated competitive performance in PTM site predictions by other researchers. In this webserver, we updated the previous framework by utilizing more advanced ensemble techniques, and providing prediction and visualization for multiple PTMs simultaneously for users to analyze potential PTM cross-talks directly. Besides prediction, users can interactively review the predicted PTM sites in the context of known PTM annotations and protein 3D structures through homology-based search. In addi- tion, the server maintains a local database providing pre-processed PTM annotations from Uniport/Swiss- Prot for users to download. This database will be up- dated every three months. The MusiteDeep server is available at https://www.musite.net. The stand-alone tools for locally using MusiteDeep are available at https://github.com/duolinwang/MusiteDeep web.
\end{abstract}

\section{Recovered Text}
W140–W146
Nucleic Acids Research, 2020, Vol. 48, Web Server issue
Published online 23 April 2020
doi: 10.1093/nar/gkaa275
MusiteDeep: a deep-learning based webserver for
protein post-translational modification site prediction
and visualization
Duolin Wang1,2, Dongpeng Liu2, Jiakang Yuchi2, Fei He1,3, Yuexu Jiang1,2, Siteng Cai2,
Jingyi Li3 and Dong Xu
1,2,*
1Christopher S. Bond Life Sciences Center, University of Missouri, Columbia, MO 65211, USA, 2Department of
Electrical Engineering and Computer Science, University of Missouri, Columbia, MO 65211, USA and 3School of
Information Science and Technology, Northeast Normal University, Changchun, Jilin 130117, China
Received February 15, 2020; Revised March 25, 2020; Editorial Decision April 07, 2020; Accepted April 08, 2020
ABSTRACT
MusiteDeep is an online resource providing a deep-
learning framework for protein post-translational
modification (PTM) site prediction and visualization.
The predictor only uses protein sequences as input
and no complex features are needed, which results
in a real-time prediction for a large number of pro-
teins. It takes less than three minutes to predict for
1000 sequences per PTM type. The output is pre-
sented at the amino acid level for the user-selected
PTM types. The framework has been benchmarked
and has demonstrated competitive performance in
PTM site predictions by other researchers. In this
webserver, we updated the previous framework by
utilizing more advanced ensemble techniques, and
providing prediction and visualization for multiple
PTMs simultaneously for users to analyze potential
PTM cross-talks directly. Besides prediction, users
can interactively review the predicted PTM sites in
the context of known PTM annotations and protein 3D
structures through homology-based search. In addi-
tion, the server maintains a local database providing
pre-processed PTM annotations from Uniport/Swiss-
Prot for users to download. This database will be up-
dated every three months. The MusiteDeep server is
available at https://www.musite.net. The stand-alone
tools for locally using MusiteDeep are available at
https://github.com/duolinwang/MusiteDeep web.
INTRODUCTION
Protein post-translational modifications (PTMs) generally
refer to the additions of various functional groups to amino
acid residues (1). PTM is a key mechanism to increase pro-
teomic diversities and plays an important role in the regu-
lation of protein functions (2). Therefore, identifying and
understanding PTMs are critical in the studies of biology
and diseases. To date, as the accumulation of PTM ex-
perimental data, dozens of bioinformatics tools have been
developed for PTM site prediction, which provide a fast
and low-cost approach in contrast to experimental meth-
ods. Many of these tools apply a machine-learning algo-
rithm and provide a prediction for a particular PTM type.
Here, we briefly review several representative tools. Musite
(3) was proposed by our group and applies a support vec-
tor machine (SVM) for phosphorylation site prediction, us-
ing the K nearest neighbor (KNN) score, disorder scores
and amino acid frequencies as features. GlycoEP (4) used
SVM for N-, O- and C-linked glycosylation site prediction.
Besides raw sequence features, GlycoEP extracted amino
acid composition (AAC) profiles, position-specific scoring
matrix (PSSM) profiles, secondary structures, and surface
accessibilities as features. GPS-PAIL (5) is a tool to pre-
dict HAT-specific lysine acetylation sites, which used their
previously proposed GPS2.2 algorithm (6) as well as fea-
tures from an amino acid substitution matrix (e.g. BLO-
SUM62) and protein-protein interaction. GPS-SUMO (7)
is another GPS algorithm-based tool to predict the sumoy-
lation sites. MePred-RF (8) applied a random forest algo-
rithm with a complex sequence-based feature selection tech-
nique to predict methylarginine and methyllysine sites. RF-
Hydroxysite (9) applied a random forest algorithm for hy-
droxylysine and hydroxyproline site prediction that com-
bined information from physicochemical, structural, evolu-
tionary and sequence-order features. Css-Palm4.0 (10) ap-
plied a clustering and scoring strategy for palmitoylation
site prediction. UbiProber (11) was designed for ubiquitin
site prediction that uses KNN, physicochemical property,
and AAC features to train an SVM-based predictor. Deep-
Phos recently applied densely connected convolutional neu-
ral network (CNN) blocks for phosphorylation site predic-
tion (12). Because the dysregulation of PTM plays impor-
*To whom correspondence should be addressed. Tel: +1 573 808 2219; Fax: +1 573 882 8318; Email: xudong@missouri.edu
C⃝The Author(s) 2020. Published by Oxford University Press on behalf of Nucleic Acids Research.
This is an Open Access article distributed under the terms of the Creative Commons Attribution Non-Commercial License
(http://creativecommons.org/licenses/by-nc/4.0/), which permits non-commercial re-use, distribution, and reproduction in any medium, provided the original work
is properly cited. For commercial re-use, please contact journals.permissions@oup.com
Downloaded from https://academic.oup.com/nar/article/48/W1/W140/5824154 by guest on 29 April 2026
Nucleic Acids Research, 2020, Vol. 48, Web Server issue W141
tant roles in the development and progression of diseases
(13–15), a number of databases were developed to annotate
existing or predicted PTM sites with SNPs or diseases, such
as dbPTM (16), PTMcode (17) and AWESOME (18). They
also provide predictions by assembling third-party tools.
For example, the recently developed AWESOME applied
20 available tools for different types of PTM site predic-
tions, and the method described in this paper was included
for phosphorylation site prediction.
Despite the availability of these tools and databases, the
number of webservers that provide a general prediction for
many different PTM types is quite low. PTM-ssMP (19) is
one, but it relies on known modification profiles, which can-
not make de-novo prediction outside the profiles. ModPred
(20) provides a general webserver for 23 different modifi-
cations. However, it was published in 2014 and applies the
logistic regression method, with a performance in need of
improvement by more advanced machine learning methods.
Most of the other existing predictive webservers were devel-
oped for a single type of PTM prediction, such as GlycoEP
and MePred-RF. All these servers limit the number of sub-
missions (mostly a single sequence), and none of them sup-
ports batch submission for any large-scale prediction, due
to calculations of more complex features, such as PSSM,
which takes significant computation time.
Here, we introduce a new webserver, MusiteDeep, to pro-
vide a general deep-learning framework for protein PTM
site prediction and visualization. The method was first in-
troduced in (21). To the best of our knowledge, MusiteDeep
was the first deep-learning method in phosphorylation (one
of the most studied PTMs) site prediction. It takes raw pro-
tein sequences as input and uses CNN (22) with a novel two-
dimensional (2D) attention mechanism. It has been widely
benchmarked by other predictors (12,23–25) and has always
ranked in first or second place. It has also been adopted
in several other services (18,26–27). Later, we upgraded the
deep-learning framework with the capsule network (28) and
extended it to predict for more PTM types, which showed
superior performance for almost all the cases with small
training samples (29). In this webserver, we updated the
previous framework by combining the two previous net-
works and utilizing more advanced ensemble techniques,
while simultaneously providing prediction and visualization
for multiple PTMs. The server provides a real-time predic-
tion for a large number of proteins by using only CPU with-
out GPU resources. Users can interactively review predicted
PTM sites in the context of known PTM annotations and
protein 3D structures through a homology-based search. In
addition, the server maintains a local database (updated ev-
ery three months) to provide pre-processed PTM annota-
tions from Uniport/Swiss-Prot for download. Comparing
with the existing web services, MusiteDeep has some ob-
vious advantages in accuracy, speed and scale, and it also
provides some unique functions for analyzing prediction re-
sults.
MATERIALS AND METHODS
Method overview
Deep-learning framework.
The framework of MusiteDeep
for protein PTM site prediction is shown in Figure 1. We
treated the PTM site prediction problem as a binary classi-
fication problem, and for each type of PTM, we trained an
independent predictor. In the training process, the frame-
work accepts raw protein sequences as input, and then 33-
length residue fragments centered at the target sites were
extracted and coded by the one-of-K coding method. We
considered the 20 common amino acids and a hyphen char-
acter ‘-’, which is used to pad the positions when the valid
fragment length is less than 33. Therefore, each position
is represented by a 21D vector, with value 1 at the index
corresponding to the amino acid or the hyphen character,
and a 0 at all other indexes; meanwhile, uncommon amino
acids are filled with 0.05. The final output of the framework
is a confidence score of the PTM prediction. In regard to
the deep-learning architecture, we used the combination of
the two previously proposed networks, i.e., MultiCNN (21)
and CapsNet (29) as shown in Figure 1 right. The Multi-
CNN model is constructed by three one-dimensional (1D)
CNN layers, a two-dimensional (2D) attention layer, and
two fully connected layers; the CapsNet is constructed by
two 1D CNN layers, a PrimaryCaps layer and a PTMCaps
layer. We trained both networks separately, and in the pre-
diction procedure, a final prediction score is calculated by
averaging the prediction scores obtained by the two inde-
pendent networks. The architecture details can be found in
the Supplementary Text S1.
Bootstrapping and weight averaging.
Because nearly all
PTM types have more negative samples than positive sam-
ples, to address the unbalanced issues during training, a
bootstrapping technique was applied. As shown in Figure 1
left, the training fragments were partitioned into N subsets,
each containing the same number of positive and negative
fragments; here, N was determined by the integer part of
the negative to positive ratio. The network was trained iter-
atively, and the number of iterations was set as N, or in prac-
tical use, an upper limit (30 by default). In each training iter-
ation, one subset was used to train the network through the
Adam stochastic optimization (30) based on mini-batches.
After training for N iterations, one classifier was obtained.
The early stopping strategy was applied in each iteration,
and when the loss of a validation set did not decrease in
some number of epochs (one forward and backward pass
over the entire subset), the training procedure for that sub-
set would stop. It has been demonstrated that an ensemble
model created by averaging the weights from a continuous
training procedure leads to wider optima and better gener-
alization (31). Therefore, during the bootstrapping proce-
dure, we applied the weight averaging strategy. We treated
each iteration trained on one subset as one training cycle
and we saved the weights generated from each iteration de-
noted as Wi, and their loss on the validation set is denoted
as Li. The final weight W of the classifier can be calculated
by the weighted average of the weights from all the iterations
as follows:
W =
N
i=1
exp (1/Li)
N
i=1 exp (1/Li)
Wi
(1)
To further improve the performance, we trained an en-
semble of 10 × N C models using a nested cross-validation.
The original training data was divided into 10 equal-sized
Downloaded from https://academic.oup.com/nar/article/48/W1/W140/5824154 by guest on 29 April 2026
W142 Nucleic Acids Research, 2020, Vol. 48, Web Server issue
Figure 1. Flowchart of the MusiteDeep framework. pos: positive fragments. neg: negative fragments.
subsets. Each time, nine subsets were used to train a model,
and the remaining subset was used as validation to moni-
tor its training process. Ten different training sets were used
during the training. In the meantime, on the same train-
ing set, we repeated the bootstrapping procedure for N C
times, and N C independent classifiers were obtained. N C
is a hyper-parameter which can be determined by users, and
N C = 2 was used in this work. The final prediction score
was the average of the scores obtained from the 10 × N C
ensemble models.
Transfer learning.
For PTM types that have small training
samples, we applied a transfer learning technique to further
improve the performance. The lower layers of CNN serve
as feature extraction layers, and the extracted features can
be generalized to a different dataset. Consider the phos-
photyrosine as an example. Because it is catalyzed by dif-
ferent kinase groups with phosphoserine or phosphothreo-
nine, we and other researchers have trained separate models
for them. However, they share some common features: For
example, the phosphorylation always happened in the dis-
order region. Therefore, in this case, the transfer learning
technique can be applied to use a larger dataset to capture
the common features and to use the smaller dataset to ex-
patiate its specified features. In particular, we trained a base
network on the phosphoserine and phosphothreonine data;
then, we used the pre-trained weights of the base network
to initialize the weights for the phosphotyrosine predictor.
Finally, we fine-tuned the weights of the predictor using the
phosphotyrosine data.
Through the combination of the two architectures and
the ensemble techniques, the overall performance has been
improved, which is shown in the Supplementary Figure S1
by comparing the previous methods on the 10-fold cross-
validation datasets provided in (29).
Performance evaluation
Because different tools used different thresholds to present
the prediction results, to avoid the effects of distinct thresh-
old values, we used the area under the ROC curves (AUC)
and the area under the precision-recall curves to evaluate
the performance. Especially, when the negative data is very
large, the precision-recall curves are more appropriate to
evaluate the performance of a predictor.
RESULTS
MusiteDeep web server
Server inputs and outputs.
The input of the server is pro-
tein sequences in the FASTA format. The server provides
two options for input, paste mode and upload mode. For
a small task with up to 10 sequences or 500 amino acids,
users can paste their sequences to the input panel, and the
job will start immediately and return the result in real time
once submitted. For a larger-scale task, users can upload
a FASTA file with up to 10 MB. In this mode, once the
job is submitted, it will be placed in a queue for process-
ing. One user can process up to five such jobs at the same
time. The job can be accessed later by the provided URL
or by checking the user’s job history. It is important to no-
tice that since we used the browser’s local storage to remem-
ber the identity of a user, the job history only lists previ-
ous jobs that were successfully submitted by the user using
the same browser; however, access through the URL has no
such restriction. All the jobs will be saved on the server for
72 hours with up to 100MB per user. MusiteDeep does not
have a complex feature calculation procedure so it can han-
dle a large number of proteins by only CPU. It takes less
than three minutes to predict for 1000 sequences per PTM
type. Users can simultaneously select multiple models for
prediction from the drop-down list. After a job is finished,
the output can be visualized for each input sequence one by
one, which can be retrieved by its sequence name or its in-
dex number. The output is shown in Figure 2. The predicted
PTMs are labeled using their abbreviations on the top of the
corresponding positions. Multiple labels are shown on top
of one position if that position is predicted to have multiple
PTMs. The highlighted colors of the predicted sites corre-
Downloaded from https://academic.oup.com/nar/article/48/W1/W140/5824154 by guest on 29 April 2026
Nucleic Acids Research, 2020, Vol. 48, Web Server issue W143
Figure 2. An example visualization of the prediction results.
spond to their prediction confidence levels. Upon hovering
the mouse on the predicted sites, the detailed information
of the prediction will be shown. A user may adjust the pre-
diction confidence threshold by using the slider to obtain
more or fewer predicted sites. Besides interactive visualiza-
tion, the results can be downloaded as plain text, with infor-
mation of protein identifier, position, residue, PTMscore,
and the predicted PTMs whose scores are higher than the
user-defined or default cut-off. An example is shown in Sup-
plementary Table S1. In the server, we also provide a REST
API for this service along with a template Python program
to demonstrate how to use the API.
View predicted PTM sites in the context of known PTM an-
notations.
MusiteDeep provides a homology-based search
against proteins in UniProtKB/Swiss-Prot, and presents
known PTM annotations at the aligned positions, as shown
in Figure 3. Protein accession identifier (ACCID), Blast se-
quence identity, and the known PTM annotations on ho-
mologous proteins are presented for each input sequence.
Upon hovering the mouse over the colored sites, the spe-
cific annotation will be shown. Proteins can be accessed in
UniProtKB/Swiss-Prot database by clicking their ACCIDs.
In the server, we also provide a REST API for this service
along with a template Python program to demonstrate how
to use the API.
View predicted PTM sites in the context of protein 3D struc-
ture.
MusiteDeep provides visualizations of the predicted
PTM sites in 3D protein structures by integrating G2S (32),
a tool to annotate genomic variants on protein structures,
and an NGL viewer (33), which is a web application for
molecular visualization. First, a query sequence with the
predicted PTM sites is searched by G2S; then, its homol-
ogous proteins that have 3D structures in RCSB PDB (34)
will be retrieved with the mapping between sequence posi-
tions and structural positions. The retrieved information is
shown in Figure 4.
The query protein and its predicted PTMs can be viewed
in the 3D structure context, as shown in Figure 5. The hover
text shows the information of the predicted site, which con-
tains its position on the query sequence, its amino acid types
at the query sequence position and at the PDB structural
position, and its predicted PTM types (in abbreviation). In
this example, multiple PTMs are viewed at the same time,
which gives the sense of 3D spatial locations of PTM sites.
Users can propose several hypotheses accordingly, such as
whether the predicted sites are physically near each other or
whether they form a potential PTM cross-talk. i.e. whether
these PTMs of multiple residues work together to determine
a particular functional outcome. In the example of Figure
5, all the predicted PTM sites locate closely on the struc-
ture, which only takes place in the loop region. Since most
PTMs tend to occur on the protein surface, we also provide
the user an option to add the molecule surface by checking
the ‘Add surface’ box.
Annotated PTM sequences from UniProtKB/Swiss-Prot.
The UniProtKB/Swiss-Prot database contains publicly
available and expertly annotated protein sequences (35), in-
cluding the PTM annotations. MusiteDeep maintains a lo-
cal database of UniProtKB/Swiss-Prot, which provides pre-
processed PTM annotations for users to download and is
updated every 3 months.
Web server implementation.
The webserver consists of a
three-layer architecture of front-end, server-end, and busi-
ness logic layers. The front-end of the server is implemented
with JavaScript libraries, React and jQuery to provide an in-
teractive user interface; the back-end is implemented with
Downloaded from https://academic.oup.com/nar/article/48/W1/W140/5824154 by guest on 29 April 2026
W144 Nucleic Acids Research, 2020, Vol. 48, Web Server issue
Figure 3. An example output of the Blast-based annotation function.
Figure 4. An example output retrieved from G2S.
Figure 5. View predicted PTM sites in the 3D structure context.
Table 1. Performance of MusiteDeep vs. ModPred and a representative method in each PTM
Area under ROC/area under precision-recall curve
PTM types
MusiteDeep
ModPred
Other
Phosphoserine/threonine
0.896/0.329
0.753/0.134
DeepPhos: 0.809/0.190
Phosphotyrosine
0.958/0.864
0.695/0.151
DeepPhos: 0.681/0.163
N-linked glycosylation
0.993/0.937
0.774/0.264
GlycoEP: 0.928/0.210
O-lined glycosylation
0.943/0.539
0.783/0.128
GlycoEP: 0.808/0.043
N6-acetyllysine
0.978/0.858
0.702/0.127
GPS-PAIL: 0.629/0.229
Methylarginine
0.941/0.844
0.770/0.130
MePred-RF: 0.681/0.152
Methyllysine
0.951/0.850
0.670/0.108
MePred-RF: 0.782/0.514
S-palmitoylation-cysteine
0.961/0.922
0.824/0.478
Css-Palm4.0: 0.735/0.465
Pyrrolidone-carboxylic-acid
0.979/0.947
0.860/0.578
-
Ubiquitination
0.804/0.279
0.584/0.091
UbiProber: 0.651/0.107
SUMOylation
0.990/0.881
0.740/0.213
GPS-SUMO: 0.706/0.357
Hydroxylysine
0.982/0.930
0.974/0.891
RF-Hydroxysite: 0.919/0.300
Hydroxyproline
0.732/0.627
0.694/0.437
RF-Hydroxysite: 0.514/0.075
Downloaded from https://academic.oup.com/nar/article/48/W1/W140/5824154 by guest on 29 April 2026
Nucleic Acids Research, 2020, Vol. 48, Web Server issue W145
KOA, a next-generation web framework for node.js; the
deep-learning framework was implemented in the business
logic layer by Python. We used mongoDB for the server’s
database. The REST APIs are implemented by KOA and
is available at the API section of the server. The stand-
alone tools used to run several services on a local machine
are available at GitHub (https://github.com/duolinwang/
MusiteDeep web) with detailed documentation.
Independent benchmarking
To demonstrate the performance of the PTM site prediction
of MusiteDeep in practical use, we compared MusiteDeep
with existing tools by using a timestamp-based dataset.
Because each type of PTM may have dozens of predic-
tive tools, to compare with all of them is very challenging.
Therefore, besides ModPred, which can cover all of our pro-
vided PTM types, we selected the most representative pub-
lic available tools based on their performance and publica-
tion citation for each PTM comparison. The training data
was constructed by extracting the protein sequences from
UniProtKB/Swiss-Prot before 2010, and the timestamp-
based testing data was constructed by the newly released
data after 2010. All the annotations generated by com-
putational predictions were removed. For each PTM, all
the residues annotated by UniProtKB/Swiss-Prot with the
same type of PTM were treated as positive sites, while the
residues with the same amino acids excluding the PTM an-
notations were regarded as the negative sites. The statistics
of the data are shown in Supplementary Table S2. The per-
formance of the whole testing data is evaluated by the area
under the ROC curves and the area under the precision-
recall curves, as shown in Table 1. We also evaluate the
performance on test subsets that have different levels of
sequence similarities to the training data. The results are
shown in Supplementary Figure S2, which shows that most
prediction performances hold well with low sequence simi-
larities.
CONCLUSIONS
In this paper, we present MusiteDeep, a web server for pro-
tein PTM site prediction and visualization. The method be-
hind the server is a combination of our previously proposed
two deep-learning models implemented with two ensemble
techniques. Since the predictor does not require the calcula-
tion of complex features, the server is capable of providing
real-time prediction and batch submission for large-scale
protein sequences. The output is presented at the amino acid
level for multiple PTMs at the same time. All the submitted
jobs will be saved in the server for 72 h with up to 100MB
for users to retrieve. Besides prediction, MusiteDeep pro-
vides facilities for users to interactively review the predicted
PTM sites in the context of known PTM annotations and
protein 3D structures. In addition, the server maintains a
local database providing pre-processed PTM annotations
from Uniport/Swiss-Prot for users to download. Compared
with the existing web services, MusiteDeep has some obvi-
ous advantages in accuracy, speed and scale. Besides the web
server, we provide user-friendly Web APIs to access several
services through Python programs and stand-alone tools,
which allow users to run MusiteDeep on local machines.
Some limitations of this work include only provid-
ing models for PTM types that have enough data in
UniProtKB/Swiss-Prot and focusing exclusively on a gen-
eral modification induced by adding functional groups to
the side chain of intermediate residues. In our future work,
we plan to extend the framework to more PTM types, in-
cluding the peptide cleavage and N-terminal PTMs, such as
N-terminal acetylation and proteolysis. We will also expand
our training data to include more databases on PTMs. The
challenge that remains is how to combine the data from dif-
ferent sources and how to control the potential errors. These
challenges will be the topics of our future studies.
DATA AVAILABILITY
MusiteDeep is available as a web server at https://
www.musite.net. The stand-alone tools to run several
MusiteDeep services on a local machine are available
in the GitHub repository (https://github.com/duolinwang/
MusiteDeep web).
SUPPLEMENTARY DATA
Supplementary Data are available at NAR Online.
FUNDING
National Institutes of Health [R35-GM126985]. Funding
for the open access charge: National Institutes of Health
[R35-GM126985].
Conflict of interest statement. None declared.
REFERENCES
1. Knorre,D.G., Kudryashova,N.V. and Godovikova,T.S. (2009)
Chemical and functional aspects of posttranslational modification of
proteins. Acta Naturae, 1, 29–51.
2. Prabakaran,S., Lippens,G., Steen,H. and Gunawardena,J. (2012)
Post-translational modification: nature’s escape from genetic
imprisonment and the basis for dynamic information encoding. Wiley
Interdiscip. Rev. Syst. Biol. Med., 4, 565–583.
3. Gao,J., Thelen,J.J., Dunker,A.K. and Xu,D. (2010) Musite, a tool for
global prediction of general and kinase-specific phosphorylation
sites. Mol. Cell. Proteomics, 9, 2586–2600.
4. Chauhan,J.S., Rao,A. and Raghava,G.P. (2013) In silico platform for
prediction of N-, O- and C-glycosites in eukaryotic protein sequences.
PLoS One, 8, e67008.
5. Deng,W., Wang,C., Zhang,Y., Xu,Y., Zhang,S., Liu,Z. and Xue,Y.
(2016) GPS-PAIL: prediction of lysine acetyltransferase-specific
modification sites from protein sequences. Sci. Rep., 6, 39787.
6. Liu,Z., Yuan,F., Ren,J., Cao,J., Zhou,Y., Yang,Q. and Xue,Y. (2012)
GPS-ARM: computational analysis of the APC/C recognition motif
by predicting D-boxes and KEN-boxes. PLoS One, 7, e34370.
7. Zhao,Q., Xie,Y., Zheng,Y., Jiang,S., Liu,W., Mu,W., Liu,Z., Zhao,Y.,
Xue,Y. and Ren,J. (2014) GPS-SUMO: a tool for the prediction of
sumoylation sites and SUMO-interaction motifs. Nucleic Acids Res.,
42, W325–W330.
8. Wei,L., Xing,P., Shi,G., Ji,Z. and Zou,Q. (2019) Fast prediction of
protein methylation sites using a sequence-based feature selection
technique. IEEE/ACM Trans. Comput. Biol. Bioinf., 16, 1264–1273.
9. Ismail,H.D., Newman,R.H. and Kc,D.B. (2016) RF-Hydroxysite: a
random forest based predictor for hydroxylation sites. Mol. Biosyst.,
12, 2427–2435.
10. Ren,J., Wen,L., Gao,X., Jin,C., Xue,Y. and Yao,X. (2008) CSS-Palm
2.0: an updated software for palmitoylation sites prediction. Protein
Eng Des Sel., 21, 639–644.
Downloaded from https://academic.oup.com/nar/article/48/W1/W140/5824154 by guest on 29 April 2026
W146 Nucleic Acids Research, 2020, Vol. 48, Web Server issue
11. Chen,X., Qiu,J.D., Shi,S.P., Suo,S.B., Huang,S.Y. and Liang,R.P.
(2013) Incorporating key position and amino acid residue features to
identify general and species-specific Ubiquitin conjugation sites.
Bioinformatics, 29, 1614–1622.
12. Luo,F., Wang,M., Liu,Y., Zhao,X.M. and Li,A. (2019) DeepPhos:
prediction of protein phosphorylation sites with deep learning.
Bioinformatics, 35, 2766–2773.
13. Santos,A.L. and Lindner,A.B. (2017) Protein posttranslational
modifications: roles in aging and age-related disease. Oxid Med Cell
Longev., 2017, 5716409.
14. Wan,L., Xu,K., Chen,Z., Tang,B. and Jiang,H. (2018) Roles of
post-translational modifications in spinocerebellar ataxias. Front.
Cell Neurosci., 12, 290.
15. Wang,Y.C., Peterson,S.E. and Loring,J.F. (2014) Protein
post-translational modifications and regulation of pluripotency in
human stem cells. Cell Res., 24, 143–160.
16. Huang,K.Y., Lee,T.Y., Kao,H.J., Ma,C.T., Lee,C.C., Lin,T.H.,
Chang,W.C. and Huang,H.D. (2019) dbPTM in 2019: exploring
disease association and cross-talk of post-translational modifications.
Nucleic Acids Res., 47, D298–D308.
17. Minguez,P., Letunic,I., Parca,L. and Bork,P. (2013) PTMcode: a
database of known and predicted functional associations between
post-translational modifications in proteins. Nucleic Acids Res., 41,
D306–D311.
18. Yang,Y., Peng,X., Ying,P., Tian,J., Li,J., Ke,J., Zhu,Y., Gong,Y.,
Zou,D., Yang,N. et al. (2019) AWESOME: a database of SNPs that
affect protein post-translational modifications. Nucleic Acids Res., 47,
D874–D880.
19. Liu,Y., Wang,M., Xi,J., Luo,F. and Li,A. (2018) PTM-ssMP: A Web
Server for Predicting Different Types of Post-translational
Modification Sites Using Novel Site-specific Modification Profile. Int.
J. Biol. Sci., 14, 946–956.
20. Pejaver,V., Hsu,W.L., Xin,F., Dunker,A.K., Uversky,V.N. and
Radivojac,P. (2014) The structural and functional signatures of
proteins that undergo multiple events of post-translational
modification. Protein Sci., 23, 1077–1093.
21. Wang,D., Zeng,S., Xu,C., Qiu,W., Liang,Y., Joshi,T. and Xu,D.
(2017) MusiteDeep: a deep-learning framework for general and
kinase-specific phosphorylation site prediction. Bioinformatics, 33,
3909–3916.
22. LeCun,Y., Bengio,Y. and Hinton,G.J.n. (2015) Deep learning. Nature,
521, 436–444.
23. Xu,Y., Song,J., Wilson,C. and Whisstock,J.C. (2018)
PhosContext2vec: a distributed representation of residue-level
sequence contexts and its application to general and kinase-specific
phosphorylation site prediction. Sci. Rep., 8, 8240.
24. Maiti,S., Hassan,A. and Mitra,P. (2020) Boosting phosphorylation
site prediction with sequence feature-based machine learning.
Proteins. 88, 284–291.
25. Fenoy,E., Izarzugaza,J.M.G., Jurtz,V., Brunak,S. and Nielsen,M.
(2019) A generic deep convolutional neural network framework for
prediction of receptor-ligand interactions-NetPhosPan: application
to kinase phosphorylation prediction. Bioinformatics, 35, 1098–1107.
26. Yu,K., Zhang,Q., Liu,Z., Zhao,Q., Zhang,X., Wang,Y., Wang,Z.X.,
Jin,Y., Li,X., Liu,Z.X. et al. (2019) qPhos: a database of protein
phosphorylation dynamics in humans. Nucleic Acids Res., 47,
D451–D458.
27. L´opez-Garc´ıa,G., Jerez,J.M., Urda,D. and Veredas,F.J. (2019) 2019
International Joint Conference on Neural Networks (IJCNN). IEEE,
1–8.
28. Sabour,S., Frosst,N. and Hinton,G.E. (2017) In: Jordan,MI,
LeCun,Y and Solla,SA (eds). Advances in Neural Information
Processing Systems. pp. 3856–3866.
29. Wang,D., Liang,Y. and Xu,D. (2019) Capsule network for protein
post-translational modification site prediction. Bioinformatics, 35,
2386–2394.
30. Kingma,D.P. and Ba,J.L. (2014) Adam: a method for stochastic
optimization. arXiv: https://arxiv.org/abs/1412.6980, 22 December
2014, pre-print: not peer reviewed.
31. Izmailov,P., Podoprikhin,D., Garipov,T., Vetrov,D. and
Wilson,A.G.J.a.p.a. (2018) Averaging weights leads to wider optima
and better generalization. arXiv: https://arxiv.org/abs/1803.05407, 14
March 2018, pre-print: not peer reviewed.
32. Wang,J., Sheridan,R., Sumer,S.O., Schultz,N., Xu,D. and Gao,J.
(2018) G2S: a web-service for annotating genomic variants on 3D
protein structures. Bioinformatics, 34, 1949–1950.
33. Rose,A.S. and Hildebrand,P.W. (2015) NGL Viewer: a web
application for molecular visualization. Nucleic Acids Res., 43,
W576–W579.
34. Burley,S.K., Berman,H.M., Bhikadiya,C., Bi,C., Chen,L., Di
Costanzo,L., Christie,C., Dalenberg,K., Duarte,J.M., Dutta,S. et al.
(2019) RCSB Protein Data Bank: biological macromolecular
structures enabling research and education in fundamental biology,
biomedicine, biotechnology and energy. Nucleic Acids Res., 47,
D464–D474.
35. Boutet,E., Lieberherr,D., Tognolli,M., Schneider,M., Bansal,P.,
Bridge,A.J., Poux,S., Bougueleret,L. and Xenarios,I. (2016)
UniProtKB/swiss-prot, the manually annotated section of the
uniprot knowledgebase: how to use the entry view. Methods Mol.
Biol., 1374, 23–54.
Downloaded from https://academic.oup.com/nar/article/48/W1/W140/5824154 by guest on 29 April 2026
\end{document}
